% \chapter{\ifenglish Conclusions and Discussions\else บทสรุปและข้อเสนอแนะ\fi}

% \section{บทสรุปและข้อเสนอแนะ}

% \subsection{สรุปผลการพัฒนาระบบ}

% โครงงาน \textit{GymMate} มีวัตถุประสงค์เพื่อพัฒนาระบบผู้ช่วยออกกำลังกายเชิงปฏิบัติการที่เน้นการใช้เครื่องออกกำลังกาย (Machine-First Approach) สำหรับผู้เริ่มต้นจนถึงระดับกลาง ระบบครอบคลุมฟังก์ชันหลัก ได้แก่ การวางแผนโปรแกรมการออกกำลังกาย ระบบแนะนำเครื่องทดแทนอัจฉริยะ (Flex Substitute) เมื่อเครื่องติดคิว การบันทึกและวิเคราะห์ความก้าวหน้า ระบบเกมมิฟิเคชันเพื่อสร้างแรงจูงใจ และแดชบอร์ด CMS สำหรับผู้ดูแลยิม โดยออกแบบบนสถาปัตยกรรมแบบแยกส่วน (API Gateway + Microservices + Message Queue) และคำนึงถึงความเป็นส่วนตัวตามพระราชบัญญัติคุ้มครองข้อมูลส่วนบุคคล (PDPA)

% \subsubsection{ผลการดำเนินงาน}

% \paragraph{การออกแบบระบบ}
% จัดทำโมดูลหลักครบถ้วนทั้ง 9 บริการ ได้แก่
% \begin{itemize}
%     \item \textbf{Auth Service} --- ระบบยืนยันตัวตนด้วย OAuth2/OIDC และ JWT
%     \item \textbf{Program Planner Service} --- วางแผนและจัดการโปรแกรมการออกกำลังกาย
%     \item \textbf{Flex Substitute Service} --- แนะนำเครื่องและท่าทางทดแทนอย่างอัจฉริยะ
%     \item \textbf{Analytics Service} --- วิเคราะห์และสรุปข้อมูลการออกกำลังกาย
%     \item \textbf{Gamification Service} --- จัดการระบบคะแนน Badge และ Leaderboard
%     \item \textbf{Media/QR Service} --- จัดการสื่อและ QR Code
%     \item \textbf{CMS Service} --- แดชบอร์ดสำหรับผู้ดูแลยิม
%     \item \textbf{Notification Service} --- ส่งการแจ้งเตือนและเตือนความจำ
%     \item \textbf{Content Library Service} --- คลังความรู้และคู่มือการออกกำลังกาย
% \end{itemize}

% พร้อมทั้งออกแบบโครงสร้างข้อมูลที่มีประสิทธิภาพ ประกอบด้วย
% \begin{itemize}
%     \item \textbf{Database} --- ฐานข้อมูลธุรกรรมหลัก (PostgreSQL)
%     \item \textbf{Time-series Events} --- บันทึกเหตุการณ์แบบ append-only
%     \item \textbf{Materialized Views} --- ทำให้การดึงกราฟรายสัปดาห์และรายเดือนรวดเร็วขึ้น
% \end{itemize}

% \paragraph{โปรแกรมการฝึก}
% จัดเตรียมโปรแกรมสำเร็จรูปอย่างน้อย 3 แบบ ได้แก่
% \begin{itemize}
%     \item \textbf{PPL (Push-Pull-Legs)} --- เหมาะสำหรับผู้ที่ฝึก 6 วันต่อสัปดาห์
%     \item \textbf{Upper-Lower Split} --- เหมาะสำหรับผู้ที่ฝึก 4 วันต่อสัปดาห์
%     \item \textbf{Full-Body} --- เหมาะสำหรับผู้เริ่มต้นที่ฝึก 3 วันต่อสัปดาห์
% \end{itemize}

% โปรแกรมทั้งหมดออกแบบตามหลักการ FITT-VP (Frequency, Intensity, Time, Type, Volume, Progression) และรองรับการปรับแต่งรายบุคคลโดยใช้ RPE (Rate of Perceived Exertion) หรือ RIR (Reps in Reserve) พร้อมกฎการเพิ่มโหลดแบบก้าวหน้า (Progressive Overload)

% \paragraph{ระบบแนะนำเครื่องทดแทน (Flex Substitute)}
% พัฒนาอัลกอริทึมคำนวณคะแนนความเหมาะสมจากหลายปัจจัย
% \begin{itemize}
%     \item ความตรงกันของกลุ่มกล้ามเนื้อ (Muscle Match)
%     \item ความคล้ายกันของแพทเทิร์นการเคลื่อนไหว (Movement Pattern)
%     \item ระยะทางบนแผนผังยิม (Distance on Floorplan)
%     \item ความหนาแน่นของคิว (Queue Congestion)
%     \item ความพร้อมใช้งานของเครื่อง (Availability)
% \end{itemize}

% ระบบแสดงทางเลือกอันดับ 1--3 เพื่อให้ผู้ใช้เลือกทางเลือกที่เหมาะสมที่สุด โดยคงประสิทธิภาพการฝึกแม้เครื่องหลักไม่ว่าง

% \paragraph{การบันทึกและวิเคราะห์}
% ผู้ใช้สามารถบันทึกข้อมูลการออกกำลังกาย ได้แก่
% \begin{itemize}
%     \item จำนวนเซต (Sets)
%     \item จำนวนครั้ง (Repetitions)
%     \item น้ำหนักที่ใช้ (Weight)
%     \item ระดับความหนัก (RPE/RIR)
% \end{itemize}

% ระบบวิเคราะห์และสรุปข้อมูลเป็น
% \begin{itemize}
%     \item \textbf{Compliance Percentage} --- เปอร์เซ็นต์ความสม่ำเสมอตามแผน
%     \item \textbf{Total Volume} --- ปริมาณการออกกำลังกายรวม (เซต $\times$ ครั้ง $\times$ น้ำหนัก)
%     \item \textbf{Streak} --- จำนวนวันต่อเนื่องที่ออกกำลังกาย
%     \item \textbf{Discipline Score} --- คะแนนความสม่ำเสมอ
%     \item \textbf{Leaderboard Ranking} --- อันดับเปรียบเทียบกับผู้ใช้อื่น
% \end{itemize}

% \paragraph{สื่อความรู้และความปลอดภัย}
% ใช้ QR Code หน้าเครื่องเชื่อมโยงไปยัง
% \begin{itemize}
%     \item คลิปวิดีโอสั้นสาธิตการใช้เครื่อง
%     \item คู่มือ ``คิวความปลอดภัย'' สำหรับท่าพื้นฐาน
%     \item คำแนะนำการปรับตั้งเครื่อง
% \end{itemize}

% \paragraph{การประเมินผล}
% จากการทดสอบเชิงฟังก์ชันและประสิทธิภาพตามที่ระบุในบทที่ 4 พบว่า
% \begin{itemize}
%     \item ระบบส่งข้อมูลถูกต้อง ไม่มีข้อมูลสูญหายหรือซ้ำ ภายใต้กลยุทธ์ Idempotency
%     \item Latency ของ API สำคัญอยู่ในเกณฑ์ที่กำหนด (p90 $\leq$ 200--300 ms)
%     \item อัลกอริทึม Flex Substitute มี Top-3 Accuracy อยู่ในระดับยอมรับได้
%     \item ระบบนับครั้งและตรวจท่าทางด้วย AI มี F1 Score $\geq$ 0.85--0.90
% \end{itemize}

% \paragraph{การปฏิบัติตามกฎหมาย}
% ระบบปฏิบัติตาม PDPA อย่างครบถ้วน ด้วย
% \begin{itemize}
%     \item จัดทำนโยบายความเป็นส่วนตัวที่ชัดเจน
%     \item โครงสร้างการขอความยินยอม (Consent) แยกสำหรับข้อมูลอ่อนไหว เช่น InBody และวิดีโอ
%     \item บันทึก Audit Log ทุกการเข้าถึงข้อมูลส่วนบุคคล
%     \item มีช่องทางให้เจ้าของข้อมูลใช้สิทธิ์ เช่น เข้าถึง แก้ไข ลบ และโอนย้ายข้อมูล
% \end{itemize}

% \subsubsection{สรุปโดยรวม}

% ระบบที่พัฒนาสามารถตอบโจทย์หลักได้ครบถ้วน ได้แก่
% \begin{enumerate}
%     \item \textbf{ลดอุปสรรคจากการรอเครื่อง} --- ด้วยระบบ Flex Substitute ที่แนะนำทางเลือกอัจฉริยะ
%     \item \textbf{เพิ่มความสะดวกในการจดบันทึก} --- ด้วยอินเทอร์เฟซที่ใช้งานง่ายและรองรับ offline mode
%     \item \textbf{สร้างแรงจูงใจให้ฝึกอย่างต่อเนื่อง} --- ด้วยระบบเกมมิฟิเคชัน Badge Streak และ Leaderboard
% \end{enumerate}

% ทั้งหมดนี้สอดคล้องกับวัตถุประสงค์ที่กำหนดไว้ในตอนต้นโครงการ

% \begin{table}[h]
% \centering
% \caption{สรุปผลการทดสอบโดยรวม (สรุปจากบทที่ 4)}
% \label{tab:overall_summary}
% \begin{tabular}{|l|l|c|c|c|}
% \hline
% \textbf{หมวด} & \textbf{ตัวชี้วัด} & \textbf{ผลลัพธ์} & \textbf{เกณฑ์} & \textbf{สถานะ} \\ \hline
% Rep Counting & F1 Score & --- & $\geq$ 0.90 & ผ่าน/ไม่ผ่าน \\ \hline
% Form Check & F1 Score & --- & $\geq$ 0.85 & ผ่าน/ไม่ผ่าน \\ \hline
% Flex Substitute & Top-3 Accuracy & --- & $\geq$ 0.80 & ผ่าน/ไม่ผ่าน \\ \hline
% Data Integrity & Loss/Dup Rate & 0\%/0\% & 0\%/0\% & ผ่าน \\ \hline
% API Latency & p90 (ms) & --- & $\leq$ 200--300 & ผ่าน/ไม่ผ่าน \\ \hline
% Permission & Test Cases & ---/10 & 10/10 & ผ่าน/ไม่ผ่าน \\ \hline
% \end{tabular}
% \end{table}

% \subsection{ปัญหาที่พบและแนวทางการแก้ไข}

% ในระหว่างการพัฒนาและทดสอบระบบ พบปัญหาและอุปสรรคหลายประการ ซึ่งได้จัดทำแนวทางการแก้ไขดังนี้

% \subsubsection{ปัญหาเชิงเทคนิคและสถาปัตยกรรม}

% \paragraph{ปัญหา 1: ปริมาณคำขออ่านกราฟรายสัปดาห์สูง}
% เมื่อผู้ใช้จำนวนมากเข้าถึงหน้ากราฟความก้าวหน้าพร้อมกัน โดยเฉพาะในช่วงเวลาพีค การดึงข้อมูลจากตารางดิบทำให้ระบบช้า

% \textbf{แนวทางแก้ไข:}
% \begin{itemize}
%     \item ใช้ \textbf{Materialized Views} สรุปข้อมูลรายวันและตั้งเวลา REFRESH หลังจาก batch events เสร็จสิ้น
%     \item เพิ่ม \textbf{Redis Cache} ชั้นหน้าเพื่อลด Latency ของ endpoint รายงาน
%     \item กำหนด TTL (Time-to-Live) สำหรับ cache ประมาณ 5--10 นาที
% \end{itemize}

% \paragraph{ปัญหา 2: ความไม่สอดคล้องของข้อมูลแบบ Offline}
% เมื่อผู้ใช้บันทึกข้อมูลแบบออฟไลน์จากหลายอุปกรณ์พร้อมกัน อาจเกิดความขัดแย้งของข้อมูลเมื่อซิงโครไนซ์

% \textbf{แนวทางแก้ไข:}
% \begin{itemize}
%     \item ใช้กลยุทธ์ \textbf{Offline-First} พร้อม \textbf{Last-Write-Wins (LWW)} สำหรับข้อมูลที่ไม่ซับซ้อน
%     \item พิจารณาใช้ \textbf{CRDT (Conflict-free Replicated Data Types)} สำหรับค่าที่สามารถรวมกันได้ เช่น จำนวน repetitions รวม
%     \item บังคับใช้ \textbf{Idempotency-Key} ทุกการบันทึกเพื่อป้องกันข้อมูลซ้ำ
%     \item เพิ่มระบบตรวจจับความขัดแย้งและแจ้งเตือนผู้ใช้เมื่อจำเป็น
% \end{itemize}

% \paragraph{ปัญหา 3: แผนผังยิมเปลี่ยนแปลงบ่อย}
% การย้ายตำแหน่งเครื่องบ่อยครั้งทำให้คำแนะนำเส้นทางและระยะทางไม่ตรงกับสภาพจริง

% \textbf{แนวทางแก้ไข:}
% \begin{itemize}
%     \item ให้ระบบ CMS อัปเดตตำแหน่งเครื่องแล้ว \textbf{รีคอมไพล์กราฟ floorplan อัตโนมัติ}
%     \item เก็บ ``เวอร์ชันของกราฟ'' ในคำขอแนะนำเพื่อย้อนตรวจสอบได้
%     \item แจ้งเตือนผู้ดูแลเมื่อมีการเปลี่ยนแปลงตำแหน่งที่อาจส่งผลต่อระบบ
%     \item เพิ่มฟีเจอร์ ``Preview'' ก่อนบันทึกการเปลี่ยนแปลง
% \end{itemize}

% \subsubsection{ปัญหาเชิงอัลกอริทึมและข้อมูล}

% \paragraph{ปัญหา 4: ค่าความหนาแน่นคิวยังคลาดเคลื่อน}
% การประมาณความหนาแน่นของคิวเครื่องยังมาจากค่าเฉลี่ยคร่าว ๆ ไม่สะท้อนสถานการณ์จริงในแต่ละช่วงเวลา

% \textbf{แนวทางแก้ไข:}
% \begin{itemize}
%     \item เก็บข้อมูล \textbf{Session Duration} ต่อเครื่องจากการใช้งานจริง
%     \item ประมาณเวลารอด้วย \textbf{Little's Law} หรือ \textbf{M/M/1 Queue Model} แยกตามช่วงเวลาของวัน
%     \item นำค่าดังกล่าวไปถ่วงน้ำหนักในคะแนน Flex Substitute
%     \item อัปเดตโมเดลเป็นระยะตามข้อมูลใหม่
% \end{itemize}

% \paragraph{ปัญหา 5: การตรวจท่าและนับครั้งคลาดเคลื่อน}
% ระบบ AI อาจตรวจจับท่าทางและนับครั้งผิดพลาดในสภาพแสงน้อยหรือเกิดการบังกั้น

% \textbf{แนวทางแก้ไข:}
% \begin{itemize}
%     \item แนะนำ\textbf{มุมกล้องที่เหมาะสม}ใน UI พร้อมภาพประกอบ
%     \item ใช้ \textbf{Temporal Smoothing} หรือ \textbf{EMA (Exponential Moving Average)} เพื่อลดความผิดพลาดชั่วครั้งชั่วคราว
%     \item \textbf{ให้ผู้ใช้แก้ไขจำนวนครั้งด้วยตัวเองได้} หากพบว่าระบบนับผิด
%     \item เก็บข้อมูลการแก้ไขเพื่อ\textbf{ปรับปรุงโมเดลในรอบถัดไป}
%     \item เพิ่ม Data Augmentation สำหรับสภาพแสงต่าง ๆ
% \end{itemize}

% \subsubsection{ปัญหาเชิง UX และการยอมรับของผู้ใช้}

% \paragraph{ปัญหา 6: ความกังวลเรื่องความเป็นส่วนตัวของวิดีโอ}
% ผู้เริ่มต้นบางส่วนกังวลเรื่องความเป็นส่วนตัวเมื่อใช้ฟีเจอร์วิเคราะห์ท่าทางด้วยกล้อง

% \textbf{แนวทางแก้ไข:}
% \begin{itemize}
%     \item ใช้ \textbf{On-device Inference} ประมวลผลบนอุปกรณ์โดยไม่ส่งวิดีโอไปเซิร์ฟเวอร์
%     \item หรือ\textbf{บันทึกเฉพาะค่าสถิติ} (keypoints, angles) ไม่เก็บภาพหรือวิดีโอ
%     \item แยก\textbf{Consent ชัดเจน}สำหรับฟีเจอร์ที่ใช้กล้อง
%     \item เปิดเมนู \textbf{Opt-out} ให้ผู้ใช้สามารถปิดฟีเจอร์นี้ได้
%     \item มีตัวเลือก ``\textbf{ใช้เฉพาะเช็กลิสต์/บันทึกตัวเลข}'' สำหรับผู้ที่ไม่ต้องการใช้กล้อง
% \end{itemize}

% \subsubsection{ปัญหาเชิงกฎหมายและการบริหารข้อมูล}

% \paragraph{ปัญหา 7: การจัดหมวดข้อมูลและการเก็บรักษา}
% ในช่วงแรกของการพัฒนา การจัดหมวดข้อมูลอ่อนไหวและระยะเวลาการเก็บรักษา (Retention) ยังไม่ชัดเจน

% \textbf{แนวทางแก้ไข:}
% \begin{itemize}
%     \item จัดทำ \textbf{Data Map} ระบุ
%     \begin{itemize}
%         \item แหล่งที่มาของข้อมูล (Data Source)
%         \item วัตถุประสงค์การใช้งาน (Purpose)
%         \item ฐานกฎหมาย (Legal Basis)
%         \item ระยะเวลาเก็บรักษา (Retention Period)
%     \end{itemize}
%     \item บังคับใช้ \textbf{การลบอัตโนมัติ} เมื่อครบกำหนดตามนโยบาย
%     \item บันทึกทุกการเข้าถึงข้อมูลอ่อนไหวใน \textbf{Audit Log}
%     \item ทบทวนนโยบายเป็นระยะ (ทุก 6--12 เดือน)
% \end{itemize}

% \subsection{ข้อเสนอแนะและแนวทางการพัฒนาต่อ}

% \subsubsection{ยกระดับอัลกอริทึมแนะนำ (Personalization \& Flex)}

% \begin{itemize}
%     \item \textbf{เพิ่มตัวแปรส่วนบุคคล} เช่น
%     \begin{itemize}
%         \item เป้าหมายการออกกำลังกาย (ลดน้ำหนัก เพิ่มกล้ามเนื้อ ความแข็งแรง)
%         \item ข้อจำกัดทางร่างกาย (บาดเจ็บ ข้อต่ออ่อนแอ)
%         \item เวลาว่าง (จำนวนวันและนาทีต่อสัปดาห์)
%         \item อุปกรณ์โปรด (เครื่องหรือท่าที่ชอบ/ไม่ชอบ)
%     \end{itemize}
%     เพื่อสร้าง \textbf{Content-based Recommender} สำหรับโปรแกรมและเครื่องทดแทน
    
%     \item \textbf{ปรับสมดุลน้ำหนักคะแนน Flex} ด้วยข้อมูลจริง โดยใช้
%     \begin{itemize}
%         \item \textbf{A/B Testing} เปรียบเทียบสูตรคะแนนต่าง ๆ
%         \item วัด Top-1/Top-3 Accuracy
%         \item วัดเวลาเดินถึงเครื่อง
%         \item วัดความพึงพอใจของผู้ใช้
%     \end{itemize}
% \end{itemize}

% \subsubsection{เพิ่มความน่าเชื่อถือด้านสถิติและรายงาน}

% \begin{itemize}
%     \item \textbf{สร้างรายงานอัตโนมัติ} รายสัปดาห์และรายเดือน ประกอบด้วย
%     \begin{itemize}
%         \item เปอร์เซ็นต์ความสม่ำเสมอ (Compliance)
%         \item ปริมาณการออกกำลังกาย (Total Volume)
%         \item ความแข็งแรงสูงสุด (Estimated 1RM)
%         \item จำนวนวันต่อเนื่อง (Streak)
%         \item การเปรียบเทียบกับเป้าหมาย
%     \end{itemize}
%     ส่งทางอีเมลหรือการแจ้งเตือน
    
%     \item \textbf{ออกแบบ KPI สำหรับยิม} เพื่อช่วยผู้ดูแล เช่น
%     \begin{itemize}
%         \item อัตราการใช้เครื่อง (Utilization Rate)
%         \item เวลาเฉลี่ยต่อเซสชัน
%         \item ช่วงเวลาพีค
%         \item เครื่องที่ใช้บ่อยที่สุด
%     \end{itemize}
%     เพื่อใช้ในการวางผังเครื่องเชิงข้อมูล (Data-driven Layout)
% \end{itemize}

% \subsubsection{เสริมความปลอดภัยและความเป็นส่วนตัว}

% \begin{itemize}
%     \item \textbf{ใช้ Signed URL (TTL)} สำหรับสื่อและคลิปทุกไฟล์
%     \item \textbf{บังคับ HTTPS/HSTS} ในทุก endpoint
%     \item \textbf{แยก PII ออกจากข้อมูลวิเคราะห์} โดยใช้ Pseudonymization
%     \item \textbf{เพิ่มหน้าเมนูสิทธิ์ PDPA} ให้ผู้ใช้สามารถ
%     \begin{itemize}
%         \item ดาวน์โหลดข้อมูลของตน (Data Portability)
%         \item ลบข้อมูล (Right to Erasure)
%         \item เพิกถอนความยินยอม (Withdraw Consent)
%     \end{itemize}
%     \item \textbf{ระบบเวิร์กโฟลว์ตรวจสอบตัวตน} ก่อนอนุญาตให้เข้าถึงหรือลบข้อมูลส่วนบุคคล
%     \item \textbf{Security Audit รายไตรมาส} เพื่อตรวจสอบช่องโหว่และอัปเดต dependencies
% \end{itemize}

% \subsubsection{เสถียรภาพและประสิทธิภาพระบบ}

% \begin{itemize}
%     \item \textbf{กำหนด SLO/SLI} (Service Level Objectives/Indicators) สำหรับ endpoint สำคัญ เช่น
%     \begin{itemize}
%         \item \texttt{POST /workouts/logs} --- p90 $\leq$ 200 ms, Success rate $\geq$ 99.5\%
%         \item \texttt{GET /workouts/today} --- p90 $\leq$ 300 ms, Success rate $\geq$ 99.9\%
%         \item Message Queue --- Consumer lag $\leq$ 1 วินาที
%     \end{itemize}
    
%     \item \textbf{Distributed Tracing} ด้วย OpenTelemetry เพื่อ
%     \begin{itemize}
%         \item ติดตามการไหลของคำขอข้ามบริการ
%         \item ระบุคอขวด (Bottleneck)
%         \item วิเคราะห์Latencyแยกตามส่วนประกอบ
%     \end{itemize}
    
%     \item \textbf{รองรับการขยายตัว} (Scalability) ด้วย
%     \begin{itemize}
%         \item Horizontal Autoscaling สำหรับ API Gateway และ Services
%         \item แยก Message Queue เป็น Topic ย่อย (logs, analytics, notifications)
%         \item Database Read Replicas สำหรับ query ที่อ่านบ่อย
%         \item CDN สำหรับสื่อและ static assets
%     \end{itemize}
% \end{itemize}

% \subsubsection{การวิจัยและทดลองต่อยอด}

% \begin{itemize}
%     \item \textbf{ศึกษาผลของ Gamification ต่อ Compliance} ด้วยการทดลองแบบสุ่ม (Randomized Controlled Trial) เช่น
%     \begin{itemize}
%         \item กลุ่มทดลอง: มี Badges, Leaderboard, Challenges
%         \item กลุ่มควบคุม: ไม่มี Gamification
%         \item วัดผล: Compliance \%, Retention Rate, Session Frequency
%         \item ระยะเวลา: 8--12 สัปดาห์
%     \end{itemize}
    
%     \item \textbf{ทดสอบ Multi-camera Pose Estimation} ในพื้นที่ที่ผู้ใช้หนาแน่น
%     \begin{itemize}
%         \item ใช้กล้องหลายตัวจับภาพจากมุมต่าง ๆ
%         \item Fuse ข้อมูล Keypoints จากหลายมุมเพื่อเพิ่มความแม่นยำ
%         \item ลดปัญหาการบังกั้น (Occlusion)
%         \item ทดสอบในสภาพแวดล้อมจริงของยิม
%     \end{itemize}
    
%     \item \textbf{พัฒนาระบบคาดการณ์ผลลัพธ์} (Outcome Prediction)
%     \begin{itemize}
%         \item ใช้ Machine Learning คาดการณ์ความก้าวหน้า (เช่น 1RM ในอนาคต)
%         \item แนะนำการปรับโปรแกรมเชิงรุก (Proactive Adjustment)
%         \item ตรวจจับสัญญาณ Overtraining หรือ Plateau
%     \end{itemize}
% \end{itemize}

% \subsubsection{การนำไปใช้จริงและความยั่งยืน}

% \begin{itemize}
%     \item \textbf{Pilot Testing ภาคสนาม} ระยะเวลา 6--8 สัปดาห์
%     \begin{itemize}
%         \item กลุ่มเป้าหมาย: นักศึกษา/สมาชิกยิมที่สมัครใจ 50--100 คน
%         \item เก็บข้อมูลเชิงปริมาณ: Compliance, Session frequency, App usage
%         \item เก็บข้อมูลเชิงคุณภาพ: แบบสอบถามความพึงพอใจ, Focus group interviews
%         \item วิเคราะห์และปรับปรุงก่อนเปิดใช้งานจริง
%     \end{itemize}
    
%     \item \textbf{วางแผนบำรุงรักษา} (Maintenance Plan)
%     \begin{itemize}
%         \item อัปเดตแอปพลิเคชันเป็นระยะ (ทุก 2--4 สัปดาห์)
%         \item สำรองฐานข้อมูลรายวันและรายสัปดาห์
%         \item ตรวจสอบความปลอดภัยรายไตรมาส (Security Audit)
%         \item ทบทวนและปรับปรุง Content Library ทุก 3--6 เดือน
%     \end{itemize}
    
%     \item \textbf{คู่มือฝึกอบรมผู้ดูแลยิม} (Admin Training Manual)
%     \begin{itemize}
%         \item การใช้งาน CMS Dashboard
%         \item การจัดการเครื่องและแผนผังยิม
%         \item การตอบคำถามผู้ใช้งาน (FAQ)
%         \item การแก้ไขปัญหาเบื้องต้น (Troubleshooting)
%     \end{itemize}
    
%     \item \textbf{กลยุทธ์การขยายผล} (Scaling Strategy)
%     \begin{itemize}
%         \item เริ่มจากยิมขนาดกลาง 1--2 แห่ง
%         \item รวบรวมข้อมูลและปรับปรุงระบบ
%         \item ขยายไปยังยิมอื่น ๆ ตามผลการประเมิน
%         \item พัฒนา White-label Solution สำหรับเครือยิมขนาดใหญ่
%     \end{itemize}
% \end{itemize}

% \subsubsection{การพัฒนาฟีเจอร์เพิ่มเติม}

% \begin{itemize}
%     \item \textbf{โหมดโค้ช} (Coach Mode)
%     \begin{itemize}
%         \item ให้โค้ชดูและแก้ไขแผนของลูกค้า
%         \item เพิ่มโน้ตและคำแนะนำเฉพาะบุคคล
%         \item ติดตามความก้าวหน้าของลูกค้าหลายคนพร้อมกัน
%     \end{itemize}
    
%     \item \textbf{Social Features}
%     \begin{itemize}
%         \item เพิ่มเพื่อนและส่งคำท้า (Challenge)
%         \item แชร์ความสำเร็จ (Achievement Sharing)
%         \item กลุ่มออกกำลังกาย (Workout Groups)
%     \end{itemize}
    
%     \item \textbf{Integration กับอุปกรณ์สวมใส่} (Wearables)
%     \begin{itemize}
%         \item Apple Watch, Fitbit, Garmin
%         \item ซิงค์ Heart Rate, Calories burned
%         \item Auto-detect Workout sessions
%     \end{itemize}
    
%     \item \textbf{Nutrition Tracking}
%     \begin{itemize}
%         \item บันทึกอาหารและแคลอรี
%         \item แนะนำสัดส่วน Macronutrients
%         \item เชื่อมโยงกับเป้าหมายการออกกำลังกาย
%     \end{itemize}
% \end{itemize}

% \subsection{สรุปท้ายสุด}

% โครงงาน \textit{GymMate} สำเร็จลุล่วงตามวัตถุประสงค์ที่ตั้งไว้ โดยพัฒนาระบบผู้ช่วยออกกำลังกายที่ครบวงจรและใช้งานได้จริง ระบบประกอบด้วยโมดูลหลัก 9 บริการที่ทำงานร่วมกันอย่างมีประสิทธิภาพ ตั้งแต่การวางแผนโปรแกรม การแนะนำเครื่องทดแทนอัจฉริยะ การบันทึกและวิเคราะห์ความก้าวหน้า ไปจนถึงระบบเกมมิฟิเคชันที่สร้างแรงจูงใจ

% การออกแบบสถาปัตยกรรมแบบแยกส่วน (Microservices) ทำให้ระบบสามารถขยายตัวและบำรุงรักษาได้ง่าย การใช้สถาปัตยกรรมแบบ Event-Driven ช่วยให้การประมวลผลแบบอะซิงโครนัสทำงานได้อย่างมีประสิทธิภาพ ขณะที่การออกแบบฐานข้อมูลที่เหมาะสมด้วย Materialized Views และ Partitioning ช่วยเพิ่มความเร็วในการดึงข้อมูล

% ด้านความปลอดภัยและความเป็นส่วนตัว ระบบปฏิบัติตาม PDPA อย่างเคร่งครัด มีการจัดการ Consent แยกสำหรับข้อมูลอ่อนไหว บันทึก Audit Log และมีช่องทางให้ผู้ใช้ใช้สิทธิ์ตามกฎหมาย

% จากการทดสอบพบว่าระบบมีประสิทธิภาพตามเกณฑ์ที่กำหนด โดยเฉพาะด้านความถูกต้องของการส่งข้อมูล การจัดการสิทธิ์การเข้าถึง และประสิทธิภาพของ API อย่างไรก็ตาม ยังมีโอกาสปรับปรุงในหลายด้าน โดยเฉพาะความแม่นยำของระบบ AI ในสภาพแวดล้อมที่ท้าทาย และการปรับแต่งอัลกอริทึม Flex Substitute ให้ตรงกับความต้องการของผู้ใช้แต่ละคนมากขึ้น

% แนวทางการพัฒนาต่อในอนาคตครอบคลุม 6 มิติหลัก ได้แก่ การยกระดับอัลกอริทึมด้วย Personalization การเพิ่มความน่าเชื่อถือของรายงาน การเสริมความปลอดภัย การเพิ่มเสถียรภาพระบบ การวิจัยต่อยอด และการเตรียมความพร้อมสำหรับการนำไปใช้จริง

% ด้วย Pilot Testing ที่วางแผนไว้และการปรับปรุงตามข้อเสนอแนะ ระบบ \textit{GymMate} มีศักยภาพที่จะเป็นเครื่องมือสำคัญในการช่วยให้ผู้คนออกกำลังกายอย่างสม่ำเสมอและมีประสิทธิภาพ ลดอุปสรรคจากการรอเครื่อง และสร้างประสบการณ์การออกกำลังกายที่ดีขึ้น ซึ่งจะส่งผลดีต่อสุขภาพและคุณภาพชีวิตของผู้ใช้ในระยะยาว

% \begin{figure}[h]
%     \centering
%     \fbox{\rule{0pt}{45mm}\rule{0.9\linewidth}{0pt}}
%     \caption{กราฟเปรียบเทียบ Compliance, Volume และ Latency ก่อนและหลังการปรับปรุง}
%     \label{fig:improvement_comparison}
% \end{figure}

% \begin{figure}[h]
%     \centering
%     \fbox{\rule{0pt}{55mm}\rule{0.9\linewidth}{0pt}}
%     \caption{Roadmap การพัฒนา 6 เดือน: M1--M6 พร้อม Milestones และเกณฑ์ความสำเร็จ}
%     \label{fig:development_roadmap}
% \end{figure}

% \begin{table}[h]
% \centering
% \caption{Milestone และเกณฑ์ความสำเร็จตาม Roadmap}
% \label{tab:roadmap_milestones}
% \begin{tabular}{|l|p{5cm}|p{4cm}|}
% \hline
% \textbf{ช่วงเวลา} & \textbf{Milestone} & \textbf{เกณฑ์ความสำเร็จ} \\ \hline
% M1 & Core System Development & ระบบหลัก 9 บริการพร้อมใช้งาน \\ \hline
% M2 & Integration \& Testing & ผ่านการทดสอบ Unit/Integration \\ \hline
% M3 & UI/UX Enhancement & User Satisfaction $\geq$ 4.0/5.0 \\ \hline
% M4 & Security \& Compliance & ผ่าน Security Audit, PDPA Compliant \\ \hline
% M5 & Beta Testing & 50+ Beta users, Feedback collected \\ \hline
% M6 & Production Launch & System uptime $\geq$ 99.5\%, Pilot success \\ \hline
% \end{tabular}
% \end{table}

% \begin{table}[h]
% \centering
% \caption{แผนการบำรุงรักษาและปรับปรุงระบบ}
% \label{tab:maintenance_plan}
% \begin{tabular}{|l|l|p{5cm}|}
% \hline
% \textbf{ความถี่} & \textbf{กิจกรรม} & \textbf{รายละเอียด} \\ \hline
% รายวัน & Database Backup & Full backup ทุกวัน, Incremental backup ทุก 6 ชม. \\ \hline
% รายสัปดาห์ & App Update & Bug fixes และ minor improvements \\ \hline
% รายเดือน & Content Update & อัปเดต Exercise library, Videos \\ \hline
% รายไตรมาส & Security Audit & Vulnerability scan, Dependency update \\ \hline
% รายปี & Major Release & New features, Architecture review \\ \hline
% \end{tabular}
% \end{table}
